% ===== PRÉAMBULE CANONIQUE — à coller VERBATIM en tête de CHAQUE .tex =====
\documentclass[11pt,a4paper]{article}
\usepackage[utf8]{inputenc}\usepackage[T1]{fontenc}\usepackage{lmodern}
\usepackage[french]{babel}
\usepackage{amsmath,amssymb,mathtools}\usepackage{icomma}
\usepackage[version=4]{mhchem}          % \ce{...} : équations & formules
\usepackage{siunitx}\sisetup{locale=FR} % \qty{}{}, \si{}, \ang{}
\usepackage{chemfig}                    % \chemfig{...} : structures développées
\usepackage{xcolor}\usepackage{colortbl}
\usepackage{tikz}\usepackage{pgfplots}\pgfplotsset{compat=1.18}
\usepackage[most]{tcolorbox}\usepackage{enumitem}\usepackage{array,booktabs}
\usepackage[a4paper,margin=2.2cm]{geometry}
\usetikzlibrary{arrows.meta,calc,angles,quotes,positioning,patterns,backgrounds,shapes.geometric}
\definecolor{primary}{RGB}{222,49,99}\definecolor{primarydark}{RGB}{150,22,55}
\definecolor{defcol}{RGB}{0,114,178}\definecolor{methcol}{RGB}{52,92,148}
\definecolor{excol}{RGB}{0,135,90}\definecolor{attncol}{RGB}{200,40,55}
\tcbset{boxbase/.style={enhanced,breakable,boxrule=0.9pt,arc=2.5pt,
  left=3mm,right=3mm,top=2.3mm,bottom=2.3mm,fonttitle=\bfseries,coltitle=white}}
% --- boîtes TUTORIEL (noms EXACTS, ne pas renommer) ---
\newtcolorbox{cle}[1][]{boxbase,colback=defcol!6,colframe=defcol,
  title={\textsc{À retenir}\ifx\relax#1\relax\relax\else\textmd{ : \itshape #1}\fi}}
\newtcolorbox{methode}[1][]{boxbase,colback=methcol!7,colframe=methcol,sharp corners,
  title={\textsc{Méthode}\ifx\relax#1\relax\relax\else\textmd{ --- \itshape #1}\fi}}
\newtcolorbox{exemplebox}[1][]{boxbase,colback=excol!5,colframe=excol,
  title={\textsc{Exemple}\ifx\relax#1\relax\relax\else\textmd{ : \itshape #1}\fi}}
\newtcolorbox{piege}[1][]{boxbase,colback=attncol!6,colframe=attncol,
  title={\textsc{Piège}\ifx\relax#1\relax\relax\else\textmd{ : \itshape #1}\fi}}
\newtcolorbox{remarque}[1][]{boxbase,colback=black!4,colframe=black!45,
  title={\textsc{Remarque}\ifx\relax#1\relax\relax\else\textmd{ : \itshape #1}\fi}}
% --- boîtes EXERCICE / CORRIGÉ (noms EXACTS, auto-comptées) ---
\newtcolorbox[auto counter]{exo}[1][]{boxbase,colback=primary!4,colframe=primary,
  title={\textsc{Exercice}~\thetcbcounter\ifx\relax#1\relax\relax\else\textmd{ --- \itshape #1}\fi}}
\newtcolorbox[auto counter]{corrige}[1][]{boxbase,colback=excol!4,colframe=excol,
  title={\textsc{Corrigé}~\thetcbcounter\ifx\relax#1\relax\relax\else\textmd{ --- \itshape #1}\fi}}
\newcommand{\R}{\mathbb{R}}\newcommand{\N}{\mathbb{N}}\newcommand{\Z}{\mathbb{Z}}\newcommand{\ds}{\displaystyle}
% (un \usepackage{fancyhdr}+en-tête est possible ; il est ignoré côté web)
% =========================================================================
\begin{document}
% THEME_TITLE: Mélanges gazeux et loi de Dalton des pressions partielles

Un récipient contient rarement un seul gaz : l'air, les gaz de combustion sont des \textbf{mélanges}.
La loi de Dalton dit comment chaque constituant participe à la pression totale.

\begin{cle}[Pression partielle et loi de Dalton]
La \textbf{pression partielle} $p_i$ d'un gaz est la pression qu'il exercerait \emph{seul} dans tout le
volume, à la même température. La \textbf{pression totale} est la \textbf{somme} des pressions
partielles :
\[ \boxed{\,p_{\text{tot}}=p_A+p_B+p_C+\dots=\sum_i p_i\,} \]
Chaque gaz se comporte comme si les autres n'existaient pas.
\end{cle}

\begin{cle}[Fraction molaire et pression partielle]
La \textbf{fraction molaire} $\chi_A=\dfrac{n_A}{n_{\text{tot}}}$ (sans unité, $\sum_i\chi_i=1$). La
pression partielle est la \emph{part} de la pression totale :
\[ \boxed{\,p_A=\chi_A\times p_{\text{tot}}\,} \]
\end{cle}

\begin{center}
\begin{tikzpicture}[scale=0.9]
  \draw[thick] (0,0) rectangle (2.1,2.4);
  \foreach \x/\y in {0.5/0.5,1.0/1.3,1.6/0.8,0.7/2.0,1.7/1.9}{\fill[defcol] (\x,\y) circle (2.6pt);}
  \node[below,font=\footnotesize] at (1.05,-0.2) {\ce{N2} : $p_A$};
  \node at (2.7,1.2) {\Large $+$};
  \draw[thick] (3.3,0) rectangle (5.4,2.4);
  \foreach \x/\y in {3.8/0.7,4.6/1.4,5.0/0.6,4.1/2.0}{\fill[primary] (\x,\y) circle (3.2pt);}
  \node[below,font=\footnotesize] at (4.35,-0.2) {\ce{O2} : $p_B$};
  \node at (6.0,1.2) {\Large $=$};
  \draw[thick] (6.6,0) rectangle (8.7,2.4);
  \foreach \x/\y in {7.0/0.5,7.5/1.3,8.1/0.8,7.2/2.0,8.2/1.9}{\fill[defcol] (\x,\y) circle (2.6pt);}
  \foreach \x/\y in {7.9/0.5,7.1/1.5,8.4/1.4,7.6/2.1}{\fill[primary] (\x,\y) circle (3.2pt);}
  \node[below,font=\footnotesize] at (7.65,-0.2) {mélange};
  \node[primarydark,above,font=\scriptsize] at (7.65,2.4) {$p_{\text{tot}}=p_A+p_B$};
\end{tikzpicture}
\end{center}

\begin{cle}[Fraction molaire $=$ fraction volumique (gaz parfaits)]
Pour des gaz parfaits, la fraction \textbf{molaire} égale la fraction \textbf{volumique} (conséquence
d'Avogadro). Ainsi « l'air contient $20\,\%$ de \ce{O2} en volume » signifie $\chi(\ce{O2})=0{,}20$,
donc $p(\ce{O2})=0{,}20\times p_{\text{atm}}$.
\end{cle}

\begin{methode}[Traiter un mélange gazeux]
\begin{enumerate}[leftmargin=1.6em,topsep=2pt,itemsep=1pt]
  \item Compter les moles de chaque gaz ($n_i=m_i/M_i$ si l'on a des masses).
  \item $n_{\text{tot}}=\sum_i n_i$, puis $\chi_i=n_i/n_{\text{tot}}$.
  \item $p_{\text{tot}}=\dfrac{n_{\text{tot}}RT}{V}$ si besoin, puis $p_i=\chi_i\,p_{\text{tot}}$.
  \item \textbf{Vérifier} : $\sum_i p_i=p_{\text{tot}}$ et $\sum_i\chi_i=1$.
\end{enumerate}
\end{methode}

\begin{exemplebox}[un mélange \ce{N2}/\ce{O2}]
$n(\ce{N2})=0{,}30$, $n(\ce{O2})=0{,}20$ mol dans $V=\qty{10.0}{\litre}$ à \qty{300}{\kelvin}. \\
$n_{\text{tot}}=0{,}50$ ; $p_{\text{tot}}=\dfrac{0{,}50\times0{,}082\times300}{10{,}0}=\qty{1.23}{atm}$. \\
$\chi(\ce{N2})=0{,}60$, $\chi(\ce{O2})=0{,}40$ ; $p(\ce{N2})=0{,}74$, $p(\ce{O2})=0{,}49$ \si{atm}
($\text{somme}=1{,}23$ \checkmark).
\end{exemplebox}

\begin{piege}[fraction molaire, pas massique]
$p_A=\chi_A\,p_{\text{tot}}$ utilise la fraction \textbf{molaire} (rapport de \emph{moles}), jamais la
fraction massique. Un gaz lourd peu abondant en moles pèse lourd mais contribue \emph{peu} à la
pression.
\end{piege}

\end{document}
